\documentclass[12pt]{article}
\usepackage{amsmath, amssymb}
\usepackage[margin=1in]{geometry}
\setlength{\parskip}{6pt}
\title{QUBO Formulation: Fragment-Based Flexible Docking Placement Selection Problem}
\date{}
\begin{document}
\maketitle

\subsection*{Fragment-Based Flexible Docking Placement Selection Problem}

\paragraph{Problem Statement.}
Let $N$ denote the number of ligand fragments. For each fragment $i \in \{0, \dots, N-1\}$, let $M_i$ be the number of candidate placements. Let $\mathcal{E}_{i,j}$ represent the protein-fragment interaction energy for the $j$-th candidate placement of fragment $i$, where $j \in \{0, \dots, M_i-1\}$. Let $\mathcal{C}$ be the set of pairwise geometric clash terms between placements belonging to different fragments, and let $P_{(i,j),(k,l)} \ge 0$ denote the clash penalty for the pair of placements $((i,j), (k,l)) \in \mathcal{C}$. Let $\mathcal{B}$ be the set of pairwise chemical bond terms, and let $R_{(i,j),(k,l)} \ge 0$ denote the bond reward for the pair $((i,j), (k,l)) \in \mathcal{B}$. The objective is to select exactly one placement per fragment to minimize the total docking score, defined as the sum of selected interaction energies, plus the sum of clash penalties between selected placements, minus the sum of bond rewards between selected placements.

\paragraph{Decision Variables.}
For each fragment $i \in \{0, \dots, N-1\}$ and each candidate placement $j \in \{0, \dots, M_i-1\}$, introduce a binary decision variable $x_{i,j} \in \{0,1\}$. The variable $x_{i,j}$ equals $1$ if and only if the $j$-th candidate placement for fragment $i$ is selected, and $0$ otherwise.

\paragraph{QUBO Derivation.}
\textbf{Step 1.} The original objective function to be minimized is:
\begin{align}
  f(\mathbf{x}) = \sum_{i=0}^{N-1} \sum_{j=0}^{M_i-1} \mathcal{E}_{i,j} x_{i,j} + \sum_{((i,j),(k,l)) \in \mathcal{C}} P_{(i,j),(k,l)} x_{i,j} x_{k,l} - \sum_{((i,j),(k,l)) \in \mathcal{B}} R_{(i,j),(k,l)} x_{i,j} x_{k,l}.
\end{align}

\textbf{Step 2.} The problem requires exactly one placement per fragment, expressed as:
\begin{align}
  \sum_{j=0}^{M_i-1} x_{i,j} = 1, \quad \forall i \in \{0, \dots, N-1\}.
\end{align}

\textbf{Step 3.} Each equality constraint is converted to a penalty term with weight $\lambda > 0$:
\begin{align}
  \Pi_i(\mathbf{x}) = \lambda \left( \sum_{j=0}^{M_i-1} x_{i,j} - 1 \right)^2.
\end{align}
Expanding the square and applying the idempotent property $x_{i,j}^2 = x_{i,j}$ for binary variables yields:
\begin{align}
  \Pi_i(\mathbf{x}) &= \lambda \left( \sum_{j=0}^{M_i-1} x_{i,j}^2 + \sum_{\substack{0 \le j, l \le M_i-1 \\ j \neq l}} x_{i,j} x_{i,l} - 2 \sum_{j=0}^{M_i-1} x_{i,j} + 1 \right) \\
  &= \lambda \left( \sum_{j=0}^{M_i-1} x_{i,j} + \sum_{\substack{0 \le j, l \le M_i-1 \\ j \neq l}} x_{i,j} x_{i,l} - 2 \sum_{j=0}^{M_i-1} x_{i,j} + 1 \right) \\
  &= \lambda \left( 1 - \sum_{j=0}^{M_i-1} x_{i,j} + \sum_{\substack{0 \le j, l \le M_i-1 \\ j \neq l}} x_{i,j} x_{i,l} \right).
\end{align}
Using the symmetry of the quadratic terms, $\sum_{j \neq l} x_{i,j} x_{i,l} = 2 \sum_{0 \le j < l \le M_i-1} x_{i,j} x_{i,l}$, the penalty contribution for fragment $i$ simplifies to:
\begin{align}
  \Pi_i(\mathbf{x}) = \lambda - \lambda \sum_{j=0}^{M_i-1} x_{i,j} + 2\lambda \sum_{0 \le j < l \le M_i-1} x_{i,j} x_{i,l}.
\end{align}

\textbf{Step 4.} The total QUBO Hamiltonian $H(\mathbf{x})$ combines the original objective and all penalty terms:
\begin{align}
  H(\mathbf{x}) = f(\mathbf{x}) + \sum_{i=0}^{N-1} \Pi_i(\mathbf{x}).
\end{align}
Substituting the expanded forms and collecting like terms gives:
\begin{align}
  H(\mathbf{x}) &= \sum_{i=0}^{N-1} \sum_{j=0}^{M_i-1} (\mathcal{E}_{i,j} - \lambda) x_{i,j} + \sum_{i=0}^{N-1} 2\lambda \sum_{0 \le j < l \le M_i-1} x_{i,j} x_{i,l} \nonumber \\
  &\quad + \sum_{((i,j),(k,l)) \in \mathcal{C}} P_{(i,j),(k,l)} x_{i,j} x_{k,l} - \sum_{((i,j),(k,l)) \in \mathcal{B}} R_{(i,j),(k,l)} x_{i,j} x_{k,l} + \sum_{i=0}^{N-1} \lambda.
\end{align}

\textbf{Step 5.} The Hamiltonian is expressed in standard QUBO form $H(\mathbf{x}) = \mathbf{x}^T Q \mathbf{x} + c$, where $\mathbf{x}$ is the vector of all binary variables. The diagonal entries of $Q$ correspond to the linear coefficients, the off-diagonal entries correspond to the quadratic coefficients, and $c$ is the constant offset:
\begin{align}
  Q_{(i,j),(i,j)} &= \mathcal{E}_{i,j} - \lambda, \quad \forall i \in \{0,\dots,N-1\}, \ j \in \{0,\dots,M_i-1\}, \\
  Q_{(i,j),(i,l)} &= 2\lambda, \quad \forall i \in \{0,\dots,N-1\}, \ 0 \le j < l \le M_i-1, \\
  Q_{(i,j),(k,l)} &= P_{(i,j),(k,l)} - R_{(i,j),(k,l)}, \quad \forall ((i,j),(k,l)) \in \mathcal{C} \cup \mathcal{B}, \\
  Q_{(i,j),(k,l)} &= 0, \quad \text{otherwise}, \\
  c &= N\lambda.
\end{align}

\paragraph{Penalty Weight Justification.}
The penalty weight $\lambda$ must be chosen sufficiently large to ensure that any infeasible assignment incurs a strictly higher energy than the optimal feasible assignment. The maximum possible reduction in the original objective function $f(\mathbf{x})$ achievable by violating a single constraint is bounded by the sum of the absolute values of all interaction parameters involving the affected variables. Specifically, let $W_i = \max_{j} |\mathcal{E}_{i,j}| + \sum_{k \neq i} \sum_{l} (|P_{(i,j),(k,l)}| + |R_{(i,j),(k,l)}|)$. Choosing $\lambda > \max_i W_i$ guarantees that the penalty term $\lambda$ dominates any feasible objective improvement, thereby enforcing constraint satisfaction at the global minimum.

\paragraph{Correctness Argument.}
Minimizing $H(\mathbf{x})$ over $\{0,1\}^{NM}$ recovers the optimal solution to the original problem. For any feasible assignment, all constraints are satisfied, so the penalty terms vanish and $H(\mathbf{x}) = f(\mathbf{x})$. For any infeasible assignment, at least one constraint is violated, incurring a penalty strictly greater than the maximum possible decrease in $f(\mathbf{x})$. Consequently, the energy of any infeasible state is strictly higher than the energy of the optimal feasible state. Thus, the global minimum of $H(\mathbf{x})$ must correspond to a feasible assignment that minimizes $f(\mathbf{x})$.

\paragraph{Illustrative Example.}
Consider a concrete instance with $N=2$ fragments and $M_0 = M_1 = 2$ candidate placements per fragment. The decision variables are indexed as $x_{0,0}, x_{0,1}, x_{1,0}, x_{1,1}$, mapped to vector indices $0, 1, 2, 3$ respectively. Let the interaction energies be $E_{00}, E_{01}, E_{10}, E_{11}$. Let the clash penalties be $P_{00,01}, P_{00,10}, P_{00,11}, P_{01,10}, P_{01,11}, P_{10,11}$, and the bond rewards be $R_{00,01}, R_{00,10}, R_{00,11}, R_{01,10}, R_{01,11}, R_{10,11}$. We set the penalty weight to $\lambda = 10000$.
Substituting these parameters into the general formulas yields the following concrete objective:
\begin{align}
  f(\mathbf{x}) &= E_{00}x_{0,0} + E_{01}x_{0,1} + E_{10}x_{1,0} + E_{11}x_{1,1} \nonumber \\
  &\quad + P_{00,01}x_{0,0}x_{0,1} + P_{00,10}x_{0,0}x_{1,0} + P_{00,11}x_{0,0}x_{1,1} \nonumber \\
  &\quad + P_{01,10}x_{0,1}x_{1,0} + P_{01,11}x_{0,1}x_{1,1} + P_{10,11}x_{1,0}x_{1,1} \nonumber \\
  &\quad - R_{00,01}x_{0,0}x_{0,1} - R_{00,10}x_{0,0}x_{1,0} - R_{00,11}x_{0,0}x_{1,1} \nonumber \\
  &\quad - R_{01,10}x_{0,1}x_{1,0} - R_{01,11}x_{0,1}x_{1,1} - R_{10,11}x_{1,0}x_{1,1}.
\end{align}
The concrete penalty terms for the two fragments are:
\begin{align}
  \Pi_0(\mathbf{x}) &= \lambda(x_{0,0} + x_{0,1} - 1)^2 = \lambda - \lambda x_{0,0} - \lambda x_{0,1} + 2\lambda x_{0,0}x_{0,1}, \\
  \Pi_1(\mathbf{x}) &= \lambda(x_{1,0} + x_{1,1} - 1)^2 = \lambda - \lambda x_{1,0} - \lambda x_{1,1} + 2\lambda x_{1,0}x_{1,1}.
\end{align}
Combining these with the objective yields the resulting Hamiltonian $H(\mathbf{x})$:
\begin{align}
  H(\mathbf{x}) &= (E_{00} - \lambda)x_{0,0} + (E_{01} - \lambda)x_{0,1} + (E_{10} - \lambda)x_{1,0} + (E_{11} - \lambda)x_{1,1} \nonumber \\
  &\quad + 2\lambda x_{0,0}x_{0,1} + 2\lambda x_{1,0}x_{1,1} \nonumber \\
  &\quad + (P_{00,01} - R_{00,01})x_{0,0}x_{0,1} + (P_{00,10} - R_{00,10})x_{0,0}x_{1,0} \nonumber \\
  &\quad + (P_{00,11} - R_{00,11})x_{0,0}x_{1,1} + (P_{01,10} - R_{01,10})x_{0,1}x_{1,0} \nonumber \\
  &\quad + (P_{01,11} - R_{01,11})x_{0,1}x_{1,1} + (P_{10,11} - R_{10,11})x_{1,0}x_{1,1} + 2\lambda.
\end{align}
The corresponding $Q$ matrix entries (using vector indices $0,1,2,3$) are:
\begin{align}
  Q_{0,0} &= E_{00} - \lambda, & Q_{1,1} &= E_{01} - \lambda, & Q_{2,2} &= E_{10} - \lambda, & Q_{3,3} &= E_{11} - \lambda, \\
  Q_{0,1} &= 2\lambda, & Q_{2,3} &= 2\lambda, & Q_{0,2} &= P_{00,10} - R_{00,10}, & Q_{0,3} &= P_{00,11} - R_{00,11}, \\
  Q_{1,2} &= P_{01,10} - R_{01,10}, & Q_{1,3} &= P_{01,11} - R_{01,11}, & Q_{2,3} &= 2\lambda + P_{10,11} - R_{10,11}.
\end{align}
The constant offset is $c = 2\lambda = 20000$. For a fully numerical instantiation, the optimal assignment $\mathbf{x}^*$ is obtained by evaluating $H(\mathbf{x})$ over all $2^4 = 16$ binary configurations, yielding the placement selection with minimum total docking score.

\end{document}